Chương này giới thiệu về hệ thống đa robot và bài toán robot phối hợp đẩy vật trên mặt phẳng 2D. Bên cạnh đó, phần tổng quan nghiên cứu chỉ ra các nghiên cứu hiện nay xoay quanh chủ đề này, các thách thức cần được giải quyết để từ đó đề ra các mục tiêu cần được thực hiện trong đồ án.
 
\section{Tổng quan về hệ thống đa robot}
Trong những năm gần đây, cùng với sự phát triển mạnh mẽ trong lĩnh vực công nghệ nói riêng và khoa học kỹ thuật nói chung, Robotics đã, đang và sẽ là một trong những ngành được nhiều nhà nghiên cứu quan tâm và định hướng phát triển. Đó là một ngành khoa học kỹ thuật cao, có tiềm năng lớn và đóng góp quan trọng cho sự phát triển công nghệ trên toàn cầu. Trong bối cảnh hiện nay, có rất nhiều công việc mang tính nguy hiểm cao, môi trường làm việc độc hại, ảnh hưởng trực tiếp đến sức khỏe và tiềm ẩn các rủi ro liên quan đến tính mạng con người. Hay một số môi trường có địa hình khó khăn cho việc di chuyển, như là vùng có hỏa hoạn, khu vực chịu ảnh hưởng của thiên tai, các môi trường nhiễm phóng xạ,... Khi mà con người không thể trực tiếp làm việc được, thì lúc này chúng ta cần đến sự hỗ trợ đến từ robot. Robot là một thiết bị có thể thực hiện các nhiệm vụ một cách tự động dưới sự điều khiển của máy tính hoặc bằng các vi mạch điện tử đã được lập trình. Với khả năng cảm nhận, xử lý thông tin và tương tác với môi trường thông qua các cơ cấu chấp hành, robot có thể phản ứng với các biến đổi của môi trường, từ đó đưa ra các lựa chọn dựa trên môi trường mà nó làm việc.

Hệ thống đa Robot (MRS-Multi Robot System) hay robot bầy đàn \cite{6304778,U.Lima2005} là một lĩnh vực nghiên cứu tương đối mới, tuy nhiên lại thu hút được nhiều sự quan tâm đến từ cộng đồng khoa học nói riêng và robotic nói chung bởi các ưu thế vượt trội so với việc sử dụng một robot đơn lẻ \cite{10.1145/3303848}. Một hệ thống đa robot, có thể đồng nhất hoặc không, bao gồm nhiều robot di động cùng làm việc, hợp tác thông qua trao đổi thông tin mà mỗi đơn vị thu thập được để cùng hoàn thành một hoặc nhiều mục tiêu chung. Do đó, việc phân chia công việc cho mỗi robot, cũng như điều khiển toàn bộ hệ thống đa robot là một chủ đề nghiên cứu quan trọng trong lĩnh vực này. Khi đem so sánh với các robot đơn lẻ, thì MRS sở hữu ưu điểm vượt trội hơn ở khả năng thực hiện song song các nhiệm vụ của chúng, từ đó cho phép chúng thực hiện nhiệm vụ hiệu quả hơn, thậm chí là thực hiện được các nhiệm vụ mà với một robot duy nhất không thể hoàn thành được. Bên cạnh đó, với việc sử dụng hệ thống phân tán, MRS sở hữu phạm vi cảm biến rộng hơn đáng kể so với một robot đơn lẻ. Các hành động phân tán cho phép chúng thực hiện các nhiệm vụ ở nhiều địa điểm khác nhau trong cùng một thời điểm. Khả năng chịu lỗi của MRS cũng cao hơn so với một robot đơn lẻ trong một số điều kiện nhất định, vì sự thất bại của một hoặc thậm chí là một vài cá thể trong nhóm có thể được bù đắp bởi các robot còn lại, chứ không phải ngay lập tức dẫn đến sự thất bại trong nhiệm vụ và dừng cả hệ thống. Hiện nay, các phương pháp tiếp cận một bầy đàn robot thường sẽ chú trọng vào 3 ưu điểm chính đó là:

\begin{itemize}
    \item \textit{Khả năng mở rộng}: Khả năng mở rộng của một bầy đàn robot được định nghĩa là khả năng hoạt động mà không phụ thuộc vào sự thay đổi về số lượng đơn vị robot tham gia, khi tăng hay giảm số lượng vẫn không ảnh hưởng đến cấu hình hệ thống. Điều này có đạt được là nhờ sử dụng cơ chế điều khiển phân tán, dựa trên thông tin cục bộ để thực hiện tác vụ, không yêu cầu thông tin toàn cục.
        \item \textit{Sự linh hoạt}: Sự linh hoạt của một bầy robot thể hiện qua khả năng thích nghi với yêu cầu mới, sự thay đổi trong môi trường xung quanh của hệ thống. Tính đơn giản trong hành vi, cơ chế phân nhiệm và tính ngẫu nhiên chính là những yếu tố mang lại cho hệ đa robot khả năng này.
    \item \textit{Độ bền}: Độ bền của hệ thống đa robot có thể được định nghĩa là mức độ mà hệ thống có thể hoạt động khi có sự cố cục bộ. Khi xuất hiện lỗi trên một đơn vị thì các đơn vị còn lại sẽ có sự bù đắp cho đơn vị bị lỗi để hệ thống tiếp tục hoạt động.
\end{itemize}

Nhìn chung, các tính chất cơ bản của một hệ thống đa robot gồm các robot di động tự trị bao gồm khả năng cảm nhận (sensing) đối với môi trường xung quanh, tính toán (computation) dự trên các thông tin thu thập được và trao đổi thông tin (information exchange) với các đơn vị khác. Các robot có khả năng tương tác và hợp tác để thực hiện nhiều nhiệm vụ cùng một lúc nhanh và linh hoạt hơn một robot đơn.
Các ứng dụng chủ yếu của hệ thống đa robot bao gồm tìm kiếm cứu hộ cứu nạn (rescue mission), giám sát (surveillance) và do thám (reconnaissance), khám phá môi trường chưa biết (exploring unknown environment), tuần tra và theo dõi (patrolling and monitoring), thu thập dữ liệu từ môi trường (data gathering)… Các nghiên cứu hiện nay về hệ thống đa robot thường xoay quanh những vấn đề chính là: nghiên cứu mô hình hệ thống dựa trên hành vi (behavior-based paradigm), kiến trúc hệ thống điều khiển (control system architecture), hợp tác trong thực thi nhiệm vụ (coordination), bài toán phân nhiệm (task allocation)\cite{doi:10.1177/0278364906063426}, bài toán xây dựng đội hình (formation synthesis). Mỗi hướng nghiên cứu đều có những bài toán riêng cần được giải quyết, trong đó có những bài toán cơ bản mang tính nền tảng của hệ thống.

Từ những ưu điểm của một hệ đa robot cũng như tính tích cực của chủ đề nghiên cứu, đồ án này sẽ tập trung nghiên cứu tập trung vào bài toán điều khiển nhiều robot phối hợp đẩy một vật thể trên mặt phẳng 2D, cũng như đề xuất một phương pháp giúp xác định vị trí trọng tâm của vật cần đẩy.